% Chapter 5: Policy Circuit Design

\chapter{Policy Circuit Design}

This chapter describes the five Circom policy circuits used in PVE-ZKP. Each
circuit encodes a distinct class of validation policy relevant to web forms.
Together, they provide a small but expressive library for the healthcare-style
registration form implemented in the prototype.

\section{Design Principles}

The policy circuits in PVE-ZKP are designed with the following principles in
mind:

\begin{itemize}[leftmargin=*]
  \item \textbf{Simplicity:} circuits should be small and easy to audit.
  \item \textbf{Composability:} circuits should be reusable across forms and
    domains with minor parameter changes.
  \item \textbf{Alignment with common policies:} the circuits target validation
    patterns frequently encountered in real-world applications.
\end{itemize}

\section{Format and Structure: \texttt{format\_check}}

The \texttt{format\_check} circuit enforces simple structural constraints on a
fixed-length character array. In the prototype, it is used to check a derived
representation of the user's name or a postcode, ensuring that each character
belongs to an allowed set and that the overall pattern matches basic rules.

Inputs include an array of ASCII codes, and constraints ensure that they fall
within allowed ranges. In more advanced versions, this circuit could encode
regular-expression-style structure.

\section{Semantic Range: \texttt{threshold\_check}}

The \texttt{threshold\_check} circuit enforces a numeric range constraint. For
example, it can require that a derived age value is at least a specified
minimum. Inputs include the numeric value and one or more threshold parameters,
and constraints ensure that comparisons hold in the finite field.

In PVE-ZKP, age is computed off-circuit from a date of birth, and the circuit
verifies that the age meets a minimum requirement (e.g., 18 years).

\section{Set Membership: \texttt{membership\_check}}

The \texttt{membership\_check} circuit tests whether a value appears in an
approved set of options, such as enumerated roles or country codes. A simple
implementation enumerates the allowed values and enforces that the input equals
one of them. More sophisticated designs could use Merkle trees or accumulators
for larger sets, but the prototype focuses on small enumerations.

\section{Hidden-Field Integrity: \texttt{integrity\_check}}

The \texttt{integrity\_check} circuit binds a hidden token to server-side state
such as a secret and a session identifier. The goal is to ensure that a hidden
form field has not been tampered with and that it corresponds to a commitment
previously issued by the server.

In a full design, the circuit would recompute a Poseidon hash (or similar)
over the secret, session identifier, and token, and enforce equality with a
public commitment input. In the current prototype, the commitment handling is
simplified and marked as a limitation to be addressed in future work.

\section{Path Boundary: \texttt{boundary\_check}}

The \texttt{boundary\_check} circuit enforces that a path-like string, such as
an upload directory prefix, stays within an allowed boundary. Inputs include a
fixed-length array representing the path and another array for the expected
prefix, and constraints ensure that the prefix matches character by character.

This circuit models a common requirement in file-handling code: ensuring that
user-controlled paths do not escape a designated directory tree. Encoding such
checks in a circuit lets the receiving system verify that this policy was
correctly applied by the validation service.

Overall, these five circuits illustrate how common validation patterns can be
translated into Circom constraints. Later work can extend this library to cover
additional policies and more complex cross-field dependencies.
